% !TeX program = pdflatex

% !TeX root = PF-04.tex

% Inventor (per express statement): Tero Keski-Valkama.

% Review route and source support in REVIEW_NOTES.md before filing.

\documentclass[pct]{sm-patent}

% Change pct to us for US Letter text. A4 drawings are separate.

\PatentTitle{Addressed predecessor binding with persistent event memory and read-time learning}

\begin{document}

\descriptionstart

\pheading{Technical field}

\pnum Persistent event-memory processors for machine-log monitoring and sequence prediction.

\pheading{Background}

\pnum A current event can match a previous event yet require the content that followed that previous event. Storing only a weighted recurrent summary can obscure this two-step association, particularly in irregularly timed or interleaved streams.

\pheading{Summary}

\pnum An event writes a key containing predecessor content into storage addressed by the event type. A current query reads the resulting successor relation. A further embodiment learns read maps from persistent source records without differentiating through their original writes.

\pheading{Brief description of the drawings}

\pnum FIG. 1 is a block diagram of an implementation. FIG. 2 shows operations of the implementation. Reference signs identify the same functional elements throughout the drawings.

\pheading{Detailed description}

\pnum In FIG. 1, 100 identifies events; 110 identifies source store; 120 identifies key writer; 130 identifies temporal read; 140 identifies forecast; 150 identifies read learning. In FIG. 2, 210 identifies retained record, 220 identifies later query and 230 identifies map update.

\pnum For event $j$ of type $x_j$, encoder 110 produces $h_j$ and predecessor message $e(x_{j-1})$. Writer 120 forms $k_j=W_k[h_j;e(x_{j-1})]$ and writes it into the slot addressed by $x_j$. The predecessor is an already revealed event; no next target is available when the write is made.

\pnum At query time $t$, reader 130 forms $q_t=W_qh_t$ and contributes $s_v=\sum_{j:x_j=v}d(t-t_j)q_t^\top k_j$ to type $v$'s race score. The decay $d$ may be exponential. An aggregated slot can store the decayed sum when the read law permits that sufficient statistic. A distinct-record store retains individual source features when later read maps must be reapplied.

\pnum For local read learning, source features are retained as numerical data after their original gradient history is detached. At the current read, present key and query maps remain differentiable. With categorical error $\delta_v=1[v=y]-p_v$, a map update contracts $\delta_v$ with the elapsed-time-weighted slot trace. This does not retroactively differentiate the original encoder write.

\pnum An episodic embodiment supplies recency, induction and null alternatives. A recency alternative votes for a past type; an induction alternative votes for an already observed successor of a matching past event; the null alternative selects a base distribution. The normalized pointer distribution is mixed with the base law by a gate in $[0,1]$. Its candidate window, document boundaries and discovery work are declared.

\pnum For anonymous interleaved machine logs, a candidate predecessor $j$ receives a type-pair duration score based on $\Delta=t-t_j$, such as $-(\log(1+\Delta/s)-\mu_{ab})^2/(2\sigma_{ab}^2)-\log\sigma_{ab}$. The race-weighted predecessor message and duration enter the event model. Fitting uses the anonymous log rather than hidden item identities. A per-step residual accumulator can issue a fault alert.

\pnum FIG. 1 shows input 100, persistent source store 110, predecessor-key writer 120, reader 130, output 140 and read-learning processor 150. FIG. 2 shows retention 210, a later query 220 and a current-map update 230. A hardware implementation keeps key records near a score engine and moves selected payload references separately from candidate keys. Capacity, candidate scans and retained source bytes are measured independently.

\pheading{Hardware implementation}

\pnum A digital implementation may use synthesizable register-transfer logic mapped to an FPGA or to an application-specific integrated circuit. Addressed state uses register banks or SRAM macros; arithmetic uses integer, fixed-point or floating-point datapaths. Ready/valid or request/acknowledge interfaces apply backpressure without accepting a message twice. Synchronous FPGA execution, locally clocked circuits and asynchronous circuits are alternative implementations, each preserving its declared causal ordering and numerical precision.

\pnum Before deployment, a reference trace can check selected addresses, values, timestamps, publication versions and output probabilities against the circuit at the declared numerical precision. Timing quantization, ties, cancellation, overflow and saturation are specified for the implementation. FPGA resource utilization, post-layout timing and measured energy are implementation-dependent quantities; no particular measured hardware result is required by this embodiment.

\claimsstart

\begin{patentclaims}

\item A computer-implemented method of processing technical events, comprising receiving an event and an already observed predecessor message; forming a key from an event state and the predecessor message; writing the key to persistent storage addressed by a type of the event; subsequently matching a current query against the storage with elapsed-time weighting; and using a resulting type score in a next-event race producing a monitoring output.

\item The method of claim 1, wherein a slot accumulates decayed keys for events having the same type.

\item The method of claim 1, wherein stored source features persist beyond a gradient-history boundary and current read maps are trained without differentiating through the original source-feature computation.

\item The method of claim 3, wherein a local update combines a race prediction error with an elapsed-time-decayed slot trace.

\item The method of claim 1, wherein a normalized mixture includes recency, observed-successor induction and null alternatives.

\item The method of claim 1, wherein candidate predecessors in an anonymous merged machine log are scored using learned event-type-pair duration parameters and a resulting predecessor message enters a machine-fault predictor.

\item An event-processing apparatus comprising persistent addressed storage, a key writer, a temporal query engine and a processor configured to perform the method of claim 1.

\item A non-transitory computer-readable storage medium storing instructions that cause a processor to perform the method of claim 1.

\end{patentclaims}

\abstractstart

An event processor writes a key formed from an event state and a preceding message into persistent storage addressed by event type. A later query reads elapsed-time-weighted keys to score successor types in a race. Source features may remain stored after their original gradient history is detached, while current read maps receive learning credit. Optional recency, induction and null alternatives provide a normalized pointer mixture. For anonymous machine logs, type-pair duration clocks select predecessor messages for fault monitoring.

\end{document}

